% Figure: schematic p-T phase diagram of a pure substance. The three
% coexistence curves (sublimation, fusion, vapourisation) meet at the
% triple point; the vapourisation curve ends at the critical point.
% Above the critical point lies the supercritical region. Curve shapes
% are schematic, not to scale.
\begin{figure}[htbp]
\centering
\begin{tikzpicture}[scale=1.0, >=latex]
  % axes
  \draw[->] (0,0) -- (8.2,0) node[right, font=\small] {$T$};
  \draw[->] (0,0) -- (0,6.0) node[above, font=\small] {$p$};
  % triple point at (2.2,1.2), critical point at (6.4,4.6)
  \coordinate (TP) at (2.2,1.2);
  \coordinate (CP) at (6.4,4.6);
  % fusion line (solid-liquid): nearly vertical, slight positive slope
  \draw[engblue, very thick] (TP) -- (2.9,5.7);
  % sublimation line (solid-vapour): from origin region up to triple point
  \draw[engred, very thick] (0.4,0.1) .. controls (1.2,0.4) and (1.8,0.8) .. (TP);
  % vapourisation line (liquid-vapour): triple point to critical point
  \draw[engorange, very thick] (TP) .. controls (3.6,2.0) and (5.2,3.4) .. (CP);
  % critical and triple point markers
  \fill[engblack] (TP) circle (1.6pt);
  \fill[engblack] (CP) circle (1.6pt);
  \node[font=\scriptsize, anchor=north east] at (TP) {triple point};
  \node[font=\scriptsize, anchor=south west] at (CP) {critical point};
  % phase-region labels
  \node[font=\small] at (1.0,3.6) {solid};
  \node[font=\small] at (4.3,4.6) {liquid};
  \node[font=\small] at (5.2,1.0) {vapour};
  \node[font=\small, align=center] at (7.2,5.2) {supercritical\\fluid};
  % faint guides marking the supercritical quadrant
  \draw[enggray, dashed] (CP) -- (6.4,6.0);
  \draw[enggray, dashed] (CP) -- (8.0,4.6);
\end{tikzpicture}
\caption[Pressure-temperature phase diagram of a pure substance]{The
pressure-temperature phase diagram of a pure substance. Three
coexistence curves separate the solid, liquid, and vapour regions: the
sublimation curve (solid-vapour), the fusion curve (solid-liquid,
nearly vertical and with a slight positive slope for most substances; a
negative slope for water), and the vapourisation curve (liquid-vapour).
The three meet at the triple point, where all three phases coexist. The
vapourisation curve terminates at the critical point; beyond it the
distinction between liquid and vapour disappears and the substance is a
supercritical fluid. The vapourisation curve is the saturation curve of
the text, whose slope obeys the Clausius-Clapeyron relation. Curve
shapes are schematic.}
\label{fig:vol04ch02:pt-phase-diagram}
\end{figure}
